\documentclass[12pt]{article}
\usepackage[utf8]{inputenc}
\usepackage[T1]{fontenc}
\usepackage{hyperref}
\title{Tamki Language Resources for K-12 Curriculum Development}
\author{Compiled by Claude Sonnet 5 (claude-sonnet-5) for Andr\'as Kornai}
\date{}
\begin{document}
\maketitle
\section{Introduction}
Tamki (also Temki; ISO 639-3 tax; Glottocode tamk1242) is an endangered East Chadic language of central Chad, in the Sokoro group (with Saba, Sokoro, Mawa and Ubi). About 500 speakers were reported in 1999. It is distinct from Tumak (tmc). The searches found no educational materials.
\section{Primary Educational Resources}
None found.
\section{Linguistic and Reference Materials}
\begin{itemize}
\item Wikipedia, Tamki language: \url{https://en.wikipedia.org/wiki/Tamki_language}
\item Wikipedia, East Chadic languages (classification context): \url{https://en.wikipedia.org/wiki/East_Chadic_languages}
\item WALS entry for Tamki: \url{https://wals.info/languoid/lect/wals_code_tmk}
\item Endangered Languages Project entry: \url{https://endangeredlanguages.com/lang/2594?hl=es}
\end{itemize}
\section{Religious and Community Resources}
None found.
\section{Digital and Audio Resources}
None found.
\section{Curriculum Development Context}
Essentially no documentation was located in the searches. Any curriculum would require new documentation work and community partnership.
\end{document}